\abstract

\begingroup
\renewcommand\keywords[1]{\indent #1}
\keywords{ИИ, ОБЛАЧНАЯ ПЛАТФОРМА, БОЛЬШАЯ ЯЗЫКОВАЯ МОДЕЛЬ, МАШИННОЕ ОБУЧЕНИЕ, БЕССЕРВЕРНЫЕ ВЫЧИСЛЕНИЯ, ХОЛОДНЫЙ СТАРТ, ИНФЕРЕНС, КЭШИРОВАНИЕ, KUBERNETES, KUBEBUILDER, CONTAINER STORAGE INTERFACE, P2P, FUSE, МИКРОСЕРВИСЫ, GOLANG, HIGHLOAD}
\endgroup

\textbf{Объект исследования} --- инфраструктура бессерверного инференса моделей машинного обучения, развернутая в среде Kubernetes с применением механизма масштабирования до нуля.

\textbf{Предмет исследования} --- методы кэширования, доставки и виртуализации доступа к файлам весов ML-моделей при холодном старте инференс-подов в Kubernetes-кластере.

\textbf{Цель работы} --- Сокращение задержки холодного старта, уменьшение нагрузку на внешние реестры моделей и повышение доступности системы при горизонтальном масштабировании путем разработки архитектуры распределенного кэширования файлов моделей машинного обучения для бессерверного инференса в Kubernetes.
Для достижения поставленной цели проведен анализ причин холодного старта и влияния роста размеров ML-моделей на инфраструктуру инференса; исследованы существующие распределенные решения для кэширования моделей --- Alluxio, JuiceFS, Dragonfly, CubeFS, Rook/Ceph; реализована специализированная система распределенного chunk-кэша, включающая агент \texttt{storage-agent}, работающий как DaemonSet на каждой worker-ноде, FUSE-клиент \texttt{storage-mounter}, предоставляющий ML-runtime стандартный файловый интерфейс, и Redis-кластер в роли реестра метаданных чанков и активных агентов; сформирована методика оценки эффективности предложенной архитектуры.

\textbf{Ключевые методы} --- трехуровневое кэширование по схеме «локальный диск --- соседняя нода --- внешний репозиторий»; FUSE-монтирование виртуальной файловой системы; P2P-доставка чанков между агентами по TCP; локальная связь через Unix Domain Socket; Redis heartbeat как реестр метаданных; атомарная запись чанков во избежание гонки и частично загруженных данных.

\textbf{Ожидаемые результаты} --- снижение количества повторных загрузок весов из внешнего репозитория, уменьшение нагрузки на HuggingFace Hub, ускорение холодного старта при наличии локального или соседнего кэша, устранение NFS как единственной точки отказа.
